\chapter{Inverse and Implicit Function Theorems, Jacobians, Homogeneity, Euler}

When a relation between variables is given implicitly, the Implicit
Function Theorem (IFT) extracts derivatives without solving for one
variable in terms of the other, exactly the situation behind
indifference curves and isoquants. The Jacobian generalises the derivative
to vector-valued maps; its determinant is the local area / volume scaling
factor, and the Inverse Function Theorem turns that determinant into a test
for whether the map can be undone near a point --- and, when it can, gives
the derivative of the inverse without ever computing the inverse.
Homogeneous functions and Euler's theorem then give a clean
structural identity that ties marginal products to total output,
underpinning constant-returns-to-scale models.

\section{Core Rules}

\begin{rulebox}{IFT, One Equation in Two Unknowns}
Suppose $F:\R^{2}\to\R$ is $C^{1}$ and $F(x_{0},y_{0})=0$ with
$F_{y}(x_{0},y_{0})\neq 0$. Then near $(x_{0},y_{0})$ there is a $C^{1}$
function $y=g(x)$ satisfying $F(x,g(x))=0$, and
\[
  \dd{y}{x} \;=\; -\frac{F_{x}}{F_{y}} .
\]
\end{rulebox}

\begin{rulebox}{What the condition $F_{y}\neq0$ rules out}
The theorem is local and its hypothesis is sharp. Where $F_{y}=0$ the
conclusion genuinely fails, and the failure is visible in the picture of the
curve $F=0$ (Chapter~5):
\begin{itemize}[leftmargin=1.6em,itemsep=0.3ex]
  \item $F_{y}=0$, $F_{x}\neq0$: a \textbf{vertical tangent}. No $y(x)$
        exists near the point, but $x=x(y)$ does, with
        $\dd{x}{y}=-F_{y}/F_{x}=0$. The theorem is symmetric in the two
        variables: at least one of the two representations survives as long
        as $\grad F\neq\mathbf{0}$.
  \item $F_{x}=F_{y}=0$: a \textbf{singular point}, typically a
        self-intersection where two branches with different slopes meet.
        Neither representation exists, and $-F_{x}/F_{y}$ degenerates to
        $0/0$.
\end{itemize}
So the working rule is: the level set of $F$ is a well-behaved curve, with a
well-defined tangent, exactly where $\grad F\neq\mathbf{0}$. This is the
same gradient condition that reappears as the \emph{constraint
qualification} in the Lagrange and Kuhn--Tucker chapters, where it rules out
precisely these kinks and crossings in the constraint set.
\end{rulebox}

\begin{rulebox}{IFT, Multivariate}
If $F(x_{1},\ldots,x_{n},y)=0$ with $F_{y}\neq 0$, then
$y=g(x_{1},\ldots,x_{n})$ locally and
\[
  \pd{y}{x_{i}} \;=\; -\frac{F_{x_{i}}}{F_{y}} .
\]
\end{rulebox}

\begin{rulebox}{MRS and MRTS via IFT}
For $U(x,y)=c$ set $F=U(x,y)-c$. Then $F_{x}=U_{x}$, $F_{y}=U_{y}$, so
\[
  \dd{y}{x}\bigg|_{U=c} = -\frac{U_{x}}{U_{y}}\quad\text{(signed)},\qquad
  \mathrm{MRS} = \frac{U_{x}}{U_{y}}\quad\text{(unsigned)} .
\]
Likewise on $Q(K,L)=Q_{0}$: $\dd{L}{K}=-Q_{K}/Q_{L}$,
$\mathrm{MRTS}=Q_{K}/Q_{L}$.
\end{rulebox}

\begin{rulebox}{Jacobian Matrix}
For $\mathbf{f}:\R^{n}\to\R^{m}$,
$\mathbf{f}(x_{1},\ldots,x_{n})=(f_{1},\ldots,f_{m})$, the \textbf{Jacobian}
is
\[
  J_{\mathbf{f}} \;=\;
  \begin{pmatrix}
    \partial f_{1}/\partial x_{1} & \cdots & \partial f_{1}/\partial x_{n}\\
    \vdots & \ddots & \vdots\\
    \partial f_{m}/\partial x_{1} & \cdots & \partial f_{m}/\partial x_{n}
  \end{pmatrix} \;\in\; \R^{m\times n} .
\]
For $m=n$, $\det J_{\mathbf{f}}$ is the local volume scaling factor. Whether
it vanishes is exactly what decides local invertibility, which is the next
three boxes.
\end{rulebox}

\begin{rulebox}{Inverse Function Theorem, one variable}
Let $f$ be continuously differentiable near $x_{0}$ and set $y_{0}=f(x_{0})$.
If
\[
  f'(x_{0})\neq0 ,
\]
then $f$ is invertible on some open interval around $x_{0}$, the inverse is
differentiable at $y_{0}$, and
\[
  \boxed{\;\bigl(f^{-1}\bigr)'(y_{0})=\frac{1}{f'(x_{0})}
        =\frac{1}{f'\bigl(f^{-1}(y_{0})\bigr)}\;}
\]
The formula is the chain rule read backwards. Differentiating the identity
$f\bigl(f^{-1}(y)\bigr)=y$ gives
$f'\bigl(f^{-1}(y)\bigr)\cdot\bigl(f^{-1}\bigr)'(y)=1$, and dividing
finishes it --- which is legitimate precisely because $f'\neq0$.

Note which point goes where: the derivative of the \emph{inverse} is
evaluated at $y_{0}$, the derivative of $f$ at $x_{0}$. Geometrically the
graph of $f^{-1}$ is the graph of $f$ mirrored in $y=x$, and mirroring turns
a slope into its reciprocal --- a vertical tangent for $f$ becoming a
horizontal one for $f^{-1}$, which is why $f'=0$ has to be excluded.

The derivatives of $\arcsin$, $\arccos$ and $\arctan$ in Chapter~30 are this
rule applied to $\sin$, $\cos$ and $\tan$; nothing else is needed to get
them.
\end{rulebox}

\begin{rulebox}{Inverse Function Theorem, several variables}
Let $\mathbf{f}:\R^{n}\to\R^{n}$ be continuously differentiable near
$\mathbf{x}_{0}$ and set $\mathbf{y}_{0}=\mathbf{f}(\mathbf{x}_{0})$. If
\[
  \det J_{\mathbf{f}}(\mathbf{x}_{0})\neq0 ,
\]
then $\mathbf{f}$ maps some open neighbourhood of $\mathbf{x}_{0}$
one-to-one onto an open neighbourhood of $\mathbf{y}_{0}$, the local inverse
is continuously differentiable, and its Jacobian is the \textbf{matrix
inverse} of the original:
\[
  \boxed{\;J_{\mathbf{f}^{-1}}(\mathbf{y}_{0})
        =\bigl[J_{\mathbf{f}}(\mathbf{x}_{0})\bigr]^{-1}\;}
\]
For $n=1$ the matrix is $1\times1$ and this is the previous box.
The condition $\det J\neq0$ is the exact analogue of $f'\neq0$: it says the
linear approximation at the point is invertible, and the theorem promises
that the map itself inherits that property nearby.

Note that the two named theorems of this chapter are the same statement
twice. Applying the Implicit Function Theorem to
$\mathbf{F}(\mathbf{x},\mathbf{y})=\mathbf{f}(\mathbf{x})-\mathbf{y}$
--- solving for $\mathbf{x}$ in terms of $\mathbf{y}$ --- has the condition
$\det J_{\mathbf{f}}\neq0$ and delivers exactly the inverse.
\end{rulebox}

\begin{rulebox}{Local, and only local}
Both theorems are \textbf{local}, and both directions of the obvious
over-reading are false.

\emph{A non-vanishing derivative everywhere does not give a global
inverse.} The map
$\mathbf{f}(x,y)=(\mathrm{e}^{x}\cos y,\ \mathrm{e}^{x}\sin y)$ has
$\det J=\mathrm{e}^{2x}>0$ at every point, yet
$\mathbf{f}(x,y)=\mathbf{f}(x,y+2\pi)$: it is invertible near each point and
nowhere invertible globally.

\emph{A vanishing derivative does not rule out an inverse.} $f(x)=x^{3}$ has
$f'(0)=0$, and is nevertheless a bijection of $\R$. What fails is not
invertibility but \emph{differentiability of the inverse}: $f^{-1}(y)=
y^{1/3}$ has a vertical tangent at $0$, exactly as the reciprocal formula
predicts.

So the condition is sufficient for a local inverse, necessary for a
\emph{differentiable} one, and silent about global behaviour.
\end{rulebox}

\begin{rulebox}{Homogeneous Functions}
$f:\R^{n}\to\R$ is \textbf{homogeneous of degree $k$} if for all $t>0$ and
all $\mathbf{x}$ in the domain,
\[
  f(t\mathbf{x}) \;=\; t^{k}f(\mathbf{x}) .
\]
Constant returns to scale (CRS) means $k=1$.
\end{rulebox}

\begin{rulebox}{Euler's Theorem}
If $f$ is differentiable and homogeneous of degree $k$ on an open cone,
then
\[
  \sum_{i=1}^{n} x_{i}\pd{f}{x_{i}} \;=\; k\,f(\mathbf{x}) .
\]
Two-variable case: $xf_{x}+yf_{y}=kf$. \textbf{Verification rule.} Always
compute the LHS $xf_{x}+yf_{y}$ and the RHS $kf$ \emph{separately}, then
check that they agree.
\end{rulebox}

\section{Worked Examples}

\begin{worked}{$F(x,y)=x^{2}y+y^{3}-10=0$ at $(1,2)$: does it define $y(x)$, and what is $y'$? Then verify Euler's theorem on $f(K,L)=K^{0.3}L^{0.7}$}
\wstep{Step 1. Check the point lies on the curve.}
$1^{2}\cdot2+2^{3}-10=2+8-10=0$. \checkmark

\wstep{Step 2. Test the hypothesis of the Implicit Function Theorem.}
The theorem needs $F_{y}\neq0$ at the point:
\[
  F_{y}=x^{2}+3y^{2},
  \qquad
  F_{y}(1,2)=1+12=13\neq0 .
\]
So there is a neighbourhood of $x=1$ on which the equation defines a unique
differentiable function $y(x)$ with $y(1)=2$.
\wtrap{The condition is on $F_{y}$ --- the derivative with respect to the
variable you are \textbf{solving for}. If you want $x$ as a function of $y$
instead, the condition is $F_{x}\neq0$. Testing the wrong partial can
license a conclusion that is false: at a point with $F_{y}=0$ the curve is
vertical and no function $y(x)$ exists there.}

\wstep{Step 3. Apply the formula.}
\[
  F_{x}=2xy,
  \qquad
  \dd{y}{x}=-\frac{F_{x}}{F_{y}}=-\frac{2xy}{x^{2}+3y^{2}},
  \qquad
  y'(1)=-\frac{4}{13}\approx-0.308 .
\]
\wtrap{The minus sign is part of the formula, and it is dropped constantly.
It comes from the chain rule applied to the identity
$F\bigl(x,y(x)\bigr)=0$: $F_{x}+F_{y}y'=0$, so $y'=-F_{x}/F_{y}$. Derive it
once and the sign stops being something to remember.}

\wstep{Step 4. Check by differentiating implicitly, from scratch.}
\[
  \frac{\mathrm{d}}{\mathrm{d}x}\bigl[x^{2}y+y^{3}-10\bigr]=0
  \;\Longrightarrow\;
  2xy+x^{2}y'+3y^{2}y'=0
  \;\Longrightarrow\;
  y'=\frac{-2xy}{x^{2}+3y^{2}} . \checkmark
\]
At $(1,2)$: $4+y'+12y'=0$, so $y'=-4/13$. The theorem's formula is not a new
technique --- it is Chapter~5 done once in general.

\wstep{Step 5. A numerical confirmation.}
Take $x=1.01$ and solve $1.0201y+y^{3}=10$ numerically: $y\approx1.996908$.
The linear prediction is $2-0.30769\times0.01=1.996923$. \checkmark\ The two
agree to four decimals over a step of $0.01$; the residual $1.5\times10^{-5}$
is the second-order term, and it should shrink by a factor of $100$ if the
step is halved twice --- which is the signature of a first-order
approximation.

\wstep{Step 6. Where this is used: comparative statics without solving.}
An equilibrium condition $F(p,\alpha)=0$ relates a price $p$ to a parameter
$\alpha$, usually with no closed-form solution for $p$. The theorem still
gives the response:
\[
  \dd{p}{\alpha}=-\frac{F_{\alpha}}{F_{p}} ,
\]
and its \emph{sign} --- which is often all that is asked --- follows from
the signs of two partials. This is the engine of every comparative-statics
result in economics.

\wstep{Step 7. Homogeneity: check the degree by scaling.}
$f$ is homogeneous of degree $r$ if $f(tK,tL)=t^{r}f(K,L)$ for all $t>0$.
For $f=K^{0.3}L^{0.7}$:
\[
  f(tK,tL)=(tK)^{0.3}(tL)^{0.7}=t^{0.3+0.7}K^{0.3}L^{0.7}=t\,f(K,L),
\]
so $r=1$: constant returns to scale. For Cobb--Douglas the degree is simply
the \emph{sum of the exponents}.
\wtrap{Scale \emph{every} argument by the same $t$, and require the identity
for all $t>0$. $f(K,L)=K^{0.3}L^{0.7}+5$ is \textbf{not} homogeneous --- the
constant does not scale --- and Euler's theorem fails for it. Homogeneity is
a strong property, not a formality.}

\wstep{Step 8. Euler's theorem, verified.}
For $f$ homogeneous of degree $r$,
\[
  K f_{K}+L f_{L}=r\,f .
\]
Here $f_{K}=0.3K^{-0.7}L^{0.7}$ and $f_{L}=0.7K^{0.3}L^{-0.3}$, so
\[
  Kf_{K}+Lf_{L}=0.3K^{0.3}L^{0.7}+0.7K^{0.3}L^{0.7}
  =(0.3+0.7)f=1\cdot f . \checkmark
\]

\wstep{Step 9. Read the economics --- this is the product-exhaustion theorem.}
Under constant returns to scale ($r=1$), paying each factor its marginal
product exhausts output exactly:
\[
  \underbrace{K\cdot MP_{K}}_{\text{capital's share}}
  +\underbrace{L\cdot MP_{L}}_{\text{labour's share}}
  =\underbrace{Q}_{\text{total output}} .
\]
Nothing is left over and nothing is missing. With increasing returns
($r>1$) the factor payments would exceed output, which is why perfect
competition and increasing returns cannot coexist --- a theorem of this
chapter, not a modelling assumption.
\end{worked}

\begin{worked}{The inverse theorem both ways: $f(x)=x^{5}+2x+1$ at $y=4$, and $\mathbf{f}(x,y)=(x^{2}-y^{2},\,2xy)$ at $(1,1)$}
\wstep{Step 1. One variable: check the hypothesis before anything else.}
\[
  f(x)=x^{5}+2x+1,
  \qquad
  f'(x)=5x^{4}+2\ \ge 2>0\quad\text{for every }x .
\]
The derivative never vanishes, so $f$ is strictly increasing and invertible
on all of $\R$, and its inverse is differentiable everywhere.

\wstep{Step 2. Locate the point --- in the \emph{domain}.}
We are asked about $y=4$, so we need the $x$ with $f(x)=4$:
\[
  x^{5}+2x+1=4
  \quad\Longrightarrow\quad
  x=1\ \ (1+2+1=4). \checkmark
\]
\wtrap{The formula needs $x_{0}=f^{-1}(y_{0})$, not $y_{0}$. Substituting
$4$ into $f'$ gives $5\cdot4^{4}+2=1282$ and the wrong answer $1/1282$.
Find the pre-image first; in an examination it is always a value you can
guess and verify, because no one can invert a quintic either.}

\wstep{Step 3. Apply the formula.}
\[
  \bigl(f^{-1}\bigr)'(4)=\frac{1}{f'(1)}=\frac{1}{5+2}=\boxed{\tfrac17}
\]
Note what was \emph{not} done: $f^{-1}$ was never computed. It cannot be ---
a general quintic has no closed-form root --- and that is exactly the point
of the theorem. It hands you a derivative of a function you cannot write
down.

\wstep{Step 4. Sanity-check it numerically.}
$f(1.01)=1.01^{5}+2.02+1\approx4.0710$, so a step of $0.0710$ in $y$ answers
to a step of $0.01$ in $x$:
\[
  \frac{\Delta x}{\Delta y}\approx\frac{0.01}{0.0710}\approx0.141
  \approx\tfrac17=0.1429 . \checkmark
\]

\wstep{Step 5. Several variables: build the Jacobian.}
For $\mathbf{f}(x,y)=(x^{2}-y^{2},\ 2xy)$,
\[
  J_{\mathbf{f}}=\begin{pmatrix}2x&-2y\\ 2y&2x\end{pmatrix},
  \qquad
  \det J_{\mathbf{f}}=4x^{2}+4y^{2}=4(x^{2}+y^{2}).
\]
This vanishes only at the origin, so the theorem applies at every other
point.

\wstep{Step 6. Evaluate at $(1,1)$ and invert the matrix.}
$\mathbf{f}(1,1)=(1-1,\,2)=(0,2)$, and
\[
  J_{\mathbf{f}}(1,1)=\begin{pmatrix}2&-2\\2&2\end{pmatrix},
  \qquad
  \det=4+4=8\neq0 .
\]
So a local inverse exists near $(0,2)$, and its Jacobian is the matrix
inverse:
\[
  J_{\mathbf{f}^{-1}}(0,2)
  =\begin{pmatrix}2&-2\\2&2\end{pmatrix}^{-1}
  =\frac18\begin{pmatrix}2&2\\-2&2\end{pmatrix}
  =\boxed{\begin{pmatrix}1/4&1/4\\-1/4&1/4\end{pmatrix}}
\]
Check by multiplying: $\begin{pmatrix}2&-2\\2&2\end{pmatrix}
\begin{pmatrix}1/4&1/4\\-1/4&1/4\end{pmatrix}
=\begin{pmatrix}1&0\\0&1\end{pmatrix}$. \checkmark
\wtrap{Invert the \emph{matrix}; do not take reciprocals entry by entry.
$\bigl[J\bigr]^{-1}$ is a matrix inverse, and the $2\times2$ rule swaps the
diagonal and negates the off-diagonal before dividing by the determinant.
Entrywise reciprocals would give $\begin{pmatrix}1/2&-1/2\\1/2&1/2
\end{pmatrix}$, which is not the inverse of anything relevant.}

\wstep{Step 7. Now show the inverse is only local.}
$\det J\neq0$ holds at every point except the origin, which invites the
conclusion that $\mathbf{f}$ is invertible on $\R^{2}\setminus\{0\}$. It is
not:
\[
  \mathbf{f}(-1,-1)=\bigl((-1)^{2}-(-1)^{2},\ 2(-1)(-1)\bigr)=(0,2)
  =\mathbf{f}(1,1).
\]
Two different points share an image, so no global inverse exists. In general
$\mathbf{f}(-x,-y)=\mathbf{f}(x,y)$: the map is two-to-one everywhere.
\wtrap{``$\det J\neq0$ everywhere'' gives a local inverse \emph{at each
point} and says nothing about fitting them together into one function. This
is the chapter's sharpest trap, and the map above is the standard
counterexample --- it is squaring, in disguise, and squaring is two-to-one.}

\wstep{Step 8. What the two halves have in common.}
\begin{center}\footnotesize
\begin{tabular}{lll}
  \hline
   & one variable & several variables\\
  \hline
  hypothesis & $f'(x_{0})\neq0$ & $\det J(\mathbf{x}_{0})\neq0$\\
  conclusion & local inverse, differentiable & the same\\
  derivative & $1/f'(x_{0})$ & $\bigl[J(\mathbf{x}_{0})\bigr]^{-1}$\\
  evaluated at & $y_{0}=f(x_{0})$ & $\mathbf{y}_{0}=\mathbf{f}(\mathbf{x}_{0})$\\
  says nothing about & global behaviour & global behaviour\\
  \hline
\end{tabular}
\end{center}
The right-hand column is the left-hand one with ``non-zero number''
replaced by ``invertible matrix'' --- which is what $\det\neq0$ means, and
why Volume~3 had to come first.
\end{worked}

\section{Practice Problems}

\begin{problems}
  \item Use IFT to find $\dd{y}{x}$ for $x^{3}+y^{3}=6xy$.
  \item Verify homogeneity and find the degree of
        $f(x,y)=x^{4}+5x^{2}y^{2}+y^{4}$.
  \item Compute the Jacobian and its determinant for
        $\mathbf{f}(x,y)=(3x-2y,\;x+4y)$.
  \item For $U(x,y)=x+2\ln y$, find the signed slope and unsigned MRS
        using IFT.
  \item Verify Euler's theorem for $f(x,y)=x^{2}+5xy-2y^{2}$ (degree $2$).
  \item Use IFT to find $\dd{y}{x}$ for $x^{2}y+xy^{3}=10$.
  \item Verify homogeneity and find the degree of
        $f(x,y)=\dfrac{x^{3}+y^{3}}{x-y}$.
  \item Compute the Jacobian and its determinant for
        $\mathbf{f}(x,y)=(x^{2}+y^{2},\;x^{2}-y^{2})$.
  \item For $U(x,y)=(x+1)(y+2)$, find the signed slope and unsigned MRS.
  \item Verify Euler's theorem for $Q(K,L)=4K^{1/4}L^{3/4}$ (degree $1$).
  \item Use IFT to find $\dd{y}{x}$ for $\mathrm{e}^{x+y}=x+y+1$.
  \item Verify homogeneity and find the degree of $f(x,y)=\sqrt{xy}$.
  \item Compute the Jacobian and its determinant for
        $\mathbf{f}(x,y)=(x\mathrm{e}^{y},\;y\mathrm{e}^{x})$.
  \item For $U(x,y)=x^{a}y^{b}$ ($a,b>0$), find the signed slope and MRS
        via IFT.
  \item Verify Euler's theorem for $Q(K,L)=K^{1/2}L^{1/2}+K+L$
        (degree $1$).
  \item Use IFT to find $\dd{y}{x}$ for $\ln(x^{2}+y)=x-y$.
  \item Verify homogeneity for
        $f(x,y,z)=\dfrac{xy+yz+zx}{x+y+z}$ and find the degree.
  \item Compute the Jacobian and its determinant for the cylindrical map
        \[
          \mathbf{f}(\rho,\theta,z)=(\rho\cos\theta,\;\rho\sin\theta,\;z).
        \]
  \item For the Leontief $U(x,y)=\min\{2x,3y\}$, discuss MRS on each
        smooth piece and at the kink.
  \item Verify Euler's theorem for the harmonic mean
        $f(x,y)=(x^{-1}+y^{-1})^{-1}$ (degree $1$).
  \item Use IFT to find $\dd{y}{x}$ for $x\cos y+y\cos x=1$.
  \item Verify homogeneity and find the degree of
        $f(x,y)=x\ln(y/x)+y\ln(x/y)$.
  \item Compute the Jacobian of
        $\mathbf{f}(x,y)=(x+y,\;xy,\;x^{2}+y^{2})$ and identify where the
        rank drops.
  \item For $Q(K,L)=5K+3L$, find the MRTS.
  \item Verify Euler's theorem for
        $Q(K,L)=\bigl(\alpha K^{\rho}+(1-\alpha)L^{\rho}\bigr)^{1/\rho}$
        (CES, degree $1$).
  \item Use IFT to find $\dd{y}{x}$ for $x^{y}=y^{x}$ with $x,y>0$.
  \item Verify whether $f(x,y)=x^{2}+xy+y$ is homogeneous; if so, find the
        degree.
  \item Compute the Jacobian and determinant of
        $\mathbf{f}(x,y)=(\mathrm{e}^{x}\cos y,\;\mathrm{e}^{x}\sin y)$.
        Interpret.
  \item For $Q(K,L)=K^{2/3}L^{1/3}$, find the MRTS via IFT.
  \item Verify Euler's theorem for $f(x,y,z)=x^{3}+y^{3}+z^{3}+3xyz$
        (degree $3$).
  \item Use IFT to find $\dd{y}{x}$ for $\sqrt{x}+\sqrt{y}+\sqrt{xy}=5$.
  \item Verify homogeneity and find the degree of
        $f(x,y)=\dfrac{x^{2}+y^{2}}{xy}$.
  % VERIFY: right-hand entries of the spherical map obscured by a printer's
  % page-turn mark in the scan; reconstructed from the standard spherical
  % parametrization, and confirmed by the worked solution below.
  \item Compute the Jacobian and determinant of the spherical map
        \[
          \mathbf{f}(\rho,\theta,\phi)=(\rho\sin\phi\cos\theta,\;
          \rho\sin\phi\sin\theta,\;\rho\cos\phi).
        \]
  \item For $U(x,y)=ax+by$ ($a,b>0$, perfect substitutes), find the MRS.
  \item Verify Euler's theorem for $f(x,y)=\dfrac{x^{3}+y^{3}}{x+y}$
        (degree $2$).
  \item Use IFT to find $\pd{z}{x}$ and $\pd{z}{y}$ for
        $F(x,y,z)=x^{2}+y^{2}+z^{2}-R^{2}=0$.
  \item Verify homogeneity of $f(x,y)=x\mathrm{e}^{y/x}$ and find the
        degree.
  \item Compute the Jacobian and determinant of
        $\mathbf{f}(u,v)=(u^{2}-v^{2},\;2uv)$.
  \item For $U(x,y)=\alpha\ln x+\beta\ln y$, find MRS via IFT.
  \item Verify Euler's theorem for $f(x,y)=\sqrt{x^{2}+y^{2}}$
        (degree $1$).
  \item Use IFT to find $\dd{y}{x}$ for $\sin(xy)+\cos(x+y)=1/2$.
  \item Decide whether $f(x,y)=x^{2}+y^{3}$ is homogeneous, with
        justification.
  \item Compute the Jacobian determinant of the polar map
        $\mathbf{f}(r,\theta)=(r\cos\theta,\;r\sin\theta)$ and connect to
        the area element.
  \item For $Q(K,L)=(K^{\rho}+L^{\rho})^{1/\rho}$ ($\rho\neq 0$), find the
        MRTS.
  \item Verify Euler's theorem for the Leontief-like
        $f(x,y)=\min\{x,y\}$ in the interior of each smooth region.
  \item Use IFT to find $\pd{y}{x_{1}}$ and $\pd{y}{x_{2}}$ at any point
        where $F(x_{1},x_{2},y)=x_{1}^{2}+x_{2}y+y^{3}-7=0$, $F_{y}\neq 0$.
  \item Verify homogeneity and find the degree of
        $f(x,y,z)=\dfrac{x^{2}y+y^{2}z+z^{2}x}{xyz}$.
  \item Compute the Jacobian determinant of
        $\mathbf{f}(x,y)=(x+\sin y,\;y+\sin x)$ and find where it
        vanishes.
  \item For $U(x,y)=(x^{\rho}+y^{\rho})^{1/\rho}$ (CES utility), find the
        MRS.
  \item Verify Euler's theorem (3-variable form) for
        $f(x,y,z)=xy+yz+zx$ (degree $2$).
  \item Use IFT to verify that on the unit circle $x^{2}+y^{2}=1$ the
        slope at $(\sqrt{3}/2,1/2)$ equals $-\sqrt{3}$.
  \item Show that $f(x,y)=(x^{1/2}+y^{1/2})^{2}$ is homogeneous of
        degree $1$, then verify Euler's theorem (LHS and RHS separately).
  \item Compute the Jacobian determinant of
        $\mathbf{f}(x,y)=(x/y,\;y/x)$.
  \item For the translog $\ln Q = \alpha_{0}+\alpha_{K}\ln K+\alpha_{L}\ln L$,
        find the MRTS in $(K,L)$ space.
  \item Verify Euler's theorem for the production function
        $Q(K,L,M)=K^{1/3}L^{1/3}M^{1/3}$ (three-factor CRS Cobb--Douglas).
  \item Compute the Jacobian and Jacobian determinant of
        $\mathbf{f}(x,y,z)=(x+y+z,\;xy+yz+zx,\;xyz)$, and explain the link
        to elementary symmetric polynomials.
  \item Show that $f(x)=x^{5}+2x+1$ is invertible on $\R$, and compute
        $\bigl(f^{-1}\bigr)'(4)$ without inverting $f$.
  \item For $f(x)=x^{3}+3x-2$, compute $\bigl(f^{-1}\bigr)'(2)$.
  \item For $f(x)=\mathrm{e}^{x}+x$, compute $\bigl(f^{-1}\bigr)'(1)$.
  \item For $f(x)=x+\ln x$ on $x>0$, compute $\bigl(f^{-1}\bigr)'(1)$.
  \item Derive $\bigl(f^{-1}\bigr)'(y)=1/f'\bigl(f^{-1}(y)\bigr)$ by
        differentiating the identity $f\bigl(f^{-1}(y)\bigr)=y$.
  \item $f(x)=x^{3}$ is a bijection of $\R$, yet $f'(0)=0$. Which
        conclusion of the Inverse Function Theorem fails at $0$, and which
        survives?
  \item Where does the Inverse Function Theorem apply to $f(x)=x^{2}$, and
        what is the largest interval containing $x=2$ on which $f$ is
        invertible?
  \item Give an interval containing $\pi/3$ on which $\sin$ is invertible,
        and compute $\bigl(\sin^{-1}\bigr)'\bigl(\tfrac{\sqrt3}{2}\bigr)$
        two ways: by the theorem, and from the known derivative of
        $\arcsin$.
  \item For $\mathbf{f}(x,y)=(x+y,\;x-y)$: compute $J$ and $\det J$, invert
        the map explicitly, and check that $J_{\mathbf{f}^{-1}}=J^{-1}$.
  \item For $\mathbf{f}(x,y)=(x^{2}-y^{2},\;2xy)$, compute $\det J$ at
        $(1,2)$ and say whether a local inverse exists there.
  \item Show that the map of the previous problem is \emph{not} globally
        injective, despite $\det J\neq0$ off the origin.
  \item For $\mathbf{f}(x,y)=(\mathrm{e}^{x}\cos y,\;\mathrm{e}^{x}\sin y)$,
        compute $\det J$ and reconcile ``non-zero everywhere'' with the fact
        that $\mathbf{f}$ is not injective.
  \item For the polar map
        $\mathbf{f}(r,\theta)=(r\cos\theta,\;r\sin\theta)$, give $\det J$
        and the set of points where the theorem gives nothing.
  \item For $\mathbf{f}(x,y)=(x+y^{3},\;y+x^{3})$, verify the hypothesis at
        the origin and give $J_{\mathbf{f}^{-1}}$ there.
  \item For $\mathbf{f}(x,y,z)=(x,\;x+y,\;x+y+z)$: compute $\det J$, write
        the inverse map explicitly, and confirm
        $J_{\mathbf{f}^{-1}}=J^{-1}$.
  \item A demand system $\mathbf{q}=\mathbf{f}(\mathbf{p})$ has
        $\det J_{\mathbf{f}}(\mathbf{p}_{0})\neq0$. What does the Inverse
        Function Theorem permit an economist to say, and what does it not?
\end{problems}

\section{Solutions}

\begin{solutions}
  \item $F=x^{3}+y^{3}-6xy$. $F_{x}=3x^{2}-6y$, $F_{y}=3y^{2}-6x$.
        \[
          \dd{y}{x} = -\frac{3x^{2}-6y}{3y^{2}-6x} = \frac{2y-x^{2}}{y^{2}-2x} .
        \]

  \item $f(tx,ty)=(tx)^{4}+5(tx)^{2}(ty)^{2}+(ty)^{4}
        =t^{4}(x^{4}+5x^{2}y^{2}+y^{4})=t^{4}f(x,y)$. \textbf{Degree}
        $k=4$.

  \item
        \[
          J=\begin{pmatrix} 3 & -2\\ 1 & 4\end{pmatrix},\qquad
          \det J=3\cdot 4-(-2)\cdot 1=14 .
        \]
        Globally invertible; areas scale by $14$.

  \item $U_{x}=1$, $U_{y}=2/y$. By IFT $\dd{y}{x}=-U_{x}/U_{y}=-y/2$.
        \textbf{MRS} $=y/2$.

  \item $f_{x}=2x+5y$, $f_{y}=5x-4y$.
        \[
          \mathrm{LHS}=x(2x+5y)+y(5x-4y)=2x^{2}+10xy-4y^{2} .
        \]
        \[
          \mathrm{RHS}=2f=2x^{2}+10xy-4y^{2} .
        \]
        $\mathrm{LHS}=\mathrm{RHS}$. \checkmark

  \item $F=x^{2}y+xy^{3}-10$. $F_{x}=2xy+y^{3}$, $F_{y}=x^{2}+3xy^{2}$.
        \[
          \dd{y}{x} = -\frac{2xy+y^{3}}{x^{2}+3xy^{2}}
            = -\frac{y(2x+y^{2})}{x(x+3y^{2})} .
        \]

  \item
        \[
          f(tx,ty)=\frac{(tx)^{3}+(ty)^{3}}{tx-ty}
            =\frac{t^{3}(x^{3}+y^{3})}{t(x-y)}=t^{2}f(x,y) .
        \]
        \textbf{Degree} $k=2$.

  \item
        \[
          J=\begin{pmatrix} 2x & 2y\\ 2x & -2y\end{pmatrix},\qquad
          \det J=(2x)(-2y)-(2y)(2x)=-8xy .
        \]
        Locally invertible whenever $x\neq 0$ and $y\neq 0$.

  \item $U_{x}=y+2$, $U_{y}=x+1$. $\dd{y}{x}=-(y+2)/(x+1)$.
        \textbf{MRS} $=(y+2)/(x+1)$.

  \item $Q_{K}=K^{-3/4}L^{3/4}$, $Q_{L}=3K^{1/4}L^{-1/4}$.
        \[
          \mathrm{LHS}=K\cdot K^{-3/4}L^{3/4}+L\cdot 3K^{1/4}L^{-1/4}
            =K^{1/4}L^{3/4}+3K^{1/4}L^{3/4}=4K^{1/4}L^{3/4} .
        \]
        \[
          \mathrm{RHS}=1\cdot Q=4K^{1/4}L^{3/4} .
        \]
        $\mathrm{LHS}=\mathrm{RHS}$. \checkmark

  \item $F=\mathrm{e}^{x+y}-x-y-1$. $F_{x}=\mathrm{e}^{x+y}-1=F_{y}$.
        \[
          \dd{y}{x} = -\frac{F_{x}}{F_{y}} = -1 .
        \]
        Valid away from $\mathrm{e}^{x+y}=1$, i.e., $x+y\neq 0$. (The
        implicit relation forces $x+y$ to be locally constant on any
        branch.)

  \item $f(tx,ty)=\sqrt{(tx)(ty)}=t\sqrt{xy}=t f(x,y)$ for $t>0$.
        \textbf{Degree} $k=1$.

  \item
        \[
          J=\begin{pmatrix} \mathrm{e}^{y} & x\mathrm{e}^{y}\\
                             y\mathrm{e}^{x} & \mathrm{e}^{x}\end{pmatrix},
          \qquad \det J=\mathrm{e}^{x+y}(1-xy) .
        \]
        Singular along $xy=1$.

  \item $U_{x}=ax^{a-1}y^{b}$, $U_{y}=bx^{a}y^{b-1}$.
        $U_{x}/U_{y}=(a/b)(y/x)$. By IFT $\dd{y}{x}=-(a/b)(y/x)$.
        \textbf{MRS} $=(a/b)(y/x)$.

  \item $Q_{K}=\tfrac12K^{-1/2}L^{1/2}+1$, $Q_{L}=\tfrac12K^{1/2}L^{-1/2}+1$.
        \[
          \mathrm{LHS}=K\bigl(\tfrac12K^{-1/2}L^{1/2}+1\bigr)
            +L\bigl(\tfrac12K^{1/2}L^{-1/2}+1\bigr)=K^{1/2}L^{1/2}+K+L .
        \]
        \[
          \mathrm{RHS}=1\cdot Q=K^{1/2}L^{1/2}+K+L .
        \]
        $\mathrm{LHS}=\mathrm{RHS}$. \checkmark

  \item $F=\ln(x^{2}+y)-x+y$. $F_{x}=\dfrac{2x}{x^{2}+y}-1$,
        $F_{y}=\dfrac{1}{x^{2}+y}+1$.
        \[
          \dd{y}{x} = -\frac{2x/(x^{2}+y)-1}{1/(x^{2}+y)+1}
            = -\frac{2x-(x^{2}+y)}{1+x^{2}+y}
            = \frac{x^{2}+y-2x}{1+x^{2}+y} .
        \]

  \item
        \[
          f(tx,ty,tz)=\frac{t^{2}(xy+yz+zx)}{t(x+y+z)}=t f(x,y,z) .
        \]
        \textbf{Degree} $k=1$.

  \item
        \[
          J=\begin{pmatrix}
            \cos\theta & -\rho\sin\theta & 0\\
            \sin\theta & \rho\cos\theta & 0\\
            0 & 0 & 1
          \end{pmatrix} .
        \]
        Expanding along the third row:
        $\det J = 1\cdot(\rho\cos^{2}\theta+\rho\sin^{2}\theta)=\rho$.
        Volume element $\mathrm{d}V=\rho\,\mathrm{d}\rho\,\mathrm{d}\theta\,\mathrm{d}z$.

  \item Leontief, non-differentiable on $\{2x=3y\}$. On $\{2x<3y\}$:
        $U=2x$, so $U_{x}=2$, $U_{y}=0$; signed slope $\dd{y}{x}$ is
        undefined (level curve is vertical), MRS is infinite: $x$ is
        the binding factor. On $\{2x>3y\}$: $U=3y$, $U_{x}=0$, $U_{y}=3$;
        signed slope $=0$, MRS $=0$: $y$ is binding. At the kink the
        two-sided derivative does not exist.

  \item $f=(x^{-1}+y^{-1})^{-1}=xy/(x+y)$. $f_{x}=y^{2}/(x+y)^{2}$,
        $f_{y}=x^{2}/(x+y)^{2}$.
        \[
          \mathrm{LHS}=x\cdot\frac{y^{2}}{(x+y)^{2}}
            +y\cdot\frac{x^{2}}{(x+y)^{2}}
            =\frac{xy(y+x)}{(x+y)^{2}}=\frac{xy}{x+y} .
        \]
        \[
          \mathrm{RHS}=1\cdot f=\frac{xy}{x+y} .
        \]
        $\mathrm{LHS}=\mathrm{RHS}$. \checkmark

  \item $F=x\cos y+y\cos x-1$. $F_{x}=\cos y-y\sin x$,
        $F_{y}=-x\sin y+\cos x$.
        \[
          \dd{y}{x} = -\frac{\cos y-y\sin x}{\cos x-x\sin y}
            = \frac{y\sin x-\cos y}{\cos x-x\sin y} .
        \]

  \item
        \[
          f(tx,ty)=(tx)\ln(ty/tx)+(ty)\ln(tx/ty)
            =tx\ln(y/x)+ty\ln(x/y)=t f(x,y)
        \]
        (the $t$ inside the log cancels). \textbf{Degree} $k=1$.

  \item
        \[
          J=\begin{pmatrix} 1 & 1\\ y & x\\ 2x & 2y\end{pmatrix}
            \in\R^{3\times 2} .
        \]
        Not square, so there is no determinant. The top $2\times 2$
        minor is $x-y$, which vanishes precisely on $x=y$; the rank
        drops to $1$ along that diagonal (the other minors also vanish
        there).

  \item Linear: $Q_{K}=5$, $Q_{L}=3$. $\dd{L}{K}=-5/3$. \textbf{MRTS}
        $=5/3$ (perfect substitutes with constant exchange rate).

  \item Let $S=\alpha K^{\rho}+(1-\alpha)L^{\rho}$, so $Q=S^{1/\rho}$.
        \[
          Q_{K}=\frac{1}{\rho}S^{1/\rho-1}\cdot\alpha\rho K^{\rho-1}
            =\alpha S^{1/\rho-1}K^{\rho-1},\qquad
          Q_{L}=(1-\alpha)S^{1/\rho-1}L^{\rho-1} .
        \]
        \[
          \mathrm{LHS}=K\alpha S^{1/\rho-1}K^{\rho-1}
            +L(1-\alpha)S^{1/\rho-1}L^{\rho-1}
            =S^{1/\rho-1}\bigl(\alpha K^{\rho}+(1-\alpha)L^{\rho}\bigr)
            =S^{1/\rho} .
        \]
        \[
          \mathrm{RHS}=1\cdot Q=S^{1/\rho} .
        \]
        $\mathrm{LHS}=\mathrm{RHS}$. \checkmark

  \item Take logs: $y\ln x-x\ln y=0$. Set $F=y\ln x-x\ln y$.
        \[
          F_{x}=\frac{y}{x}-\ln y,\qquad F_{y}=\ln x-\frac{x}{y} .
        \]
        \[
          \dd{y}{x} = -\frac{y/x-\ln y}{\ln x-x/y}
            = \frac{\ln y-y/x}{\ln x-x/y}
            = \frac{y(x\ln y-y)}{x(y\ln x-x)} .
        \]

  \item Test: $f(tx,ty)=(tx)^{2}+(tx)(ty)+(ty)=t^{2}x^{2}+t^{2}xy+ty$.
        This cannot be written as $t^{k}(x^{2}+xy+y)$ for a single $k$:
        the first two terms scale by $t^{2}$, the last by $t$.
        \textbf{Not homogeneous.}

  \item $f_{1}=\mathrm{e}^{x}\cos y$, $f_{2}=\mathrm{e}^{x}\sin y$.
        \[
          J=\begin{pmatrix} \mathrm{e}^{x}\cos y & -\mathrm{e}^{x}\sin y\\
                             \mathrm{e}^{x}\sin y & \mathrm{e}^{x}\cos y
            \end{pmatrix},\qquad
          \det J=\mathrm{e}^{2x}(\cos^{2}y+\sin^{2}y)=\mathrm{e}^{2x} .
        \]
        Always strictly positive: the map is locally invertible
        everywhere. (It is the complex exponential $w=\mathrm{e}^{x+iy}$;
        conformal, orientation-preserving.)

  \item $Q_{K}=\tfrac23K^{-1/3}L^{1/3}$, $Q_{L}=\tfrac13K^{2/3}L^{-2/3}$.
        \[
          \frac{Q_{K}}{Q_{L}}=2\,\frac{L}{K} .
        \]
        $\dd{L}{K}=-2L/K$. \textbf{MRTS} $=2L/K$.

  \item $f_{x}=3x^{2}+3yz$, $f_{y}=3y^{2}+3xz$, $f_{z}=3z^{2}+3xy$.
        \[
          \mathrm{LHS}=x(3x^{2}+3yz)+y(3y^{2}+3xz)+z(3z^{2}+3xy)
            =3(x^{3}+y^{3}+z^{3})+9xyz=3f .
        \]
        \[
          \mathrm{RHS}=3f .
        \]
        $\mathrm{LHS}=\mathrm{RHS}$. \checkmark

  \item $F=\sqrt{x}+\sqrt{y}+\sqrt{xy}-5$. Compute:
        \[
          F_{x}=\frac{1}{2\sqrt{x}}+\frac{y}{2\sqrt{xy}}
            =\frac{1}{2\sqrt{x}}+\frac{1}{2}\sqrt{\frac{y}{x}}
            =\frac{1+\sqrt{y}}{2\sqrt{x}} .
        \]
        By symmetry, $F_{y}=\dfrac{1+\sqrt{x}}{2\sqrt{y}}$.
        \[
          \dd{y}{x} = -\frac{F_{x}}{F_{y}}
            = -\frac{\sqrt{y}(1+\sqrt{y})}{\sqrt{x}(1+\sqrt{x})} .
        \]

  \item
        \[
          f(tx,ty)=\frac{t^{2}x^{2}+t^{2}y^{2}}{(tx)(ty)}
            =\frac{t^{2}(x^{2}+y^{2})}{t^{2}xy}=f(x,y) .
        \]
        \textbf{Degree} $k=0$.

  \item
        \[
          J=\begin{pmatrix}
            \sin\phi\cos\theta & -\rho\sin\phi\sin\theta & \rho\cos\phi\cos\theta\\
            \sin\phi\sin\theta & \rho\sin\phi\cos\theta & \rho\cos\phi\sin\theta\\
            \cos\phi & 0 & -\rho\sin\phi
          \end{pmatrix} .
        \]
        Expand along the third row:
        \[
          \det J = \cos\phi\Bigl[(-\rho\sin\phi\sin\theta)(\rho\cos\phi\sin\theta)
            -(\rho\cos\phi\cos\theta)(\rho\sin\phi\cos\theta)\Bigr]
        \]
        \[
          \qquad{}+(-\rho\sin\phi)\Bigl[(\sin\phi\cos\theta)(\rho\sin\phi\cos\theta)
            -(-\rho\sin\phi\sin\theta)(\sin\phi\sin\theta)\Bigr] .
        \]
        First bracket:
        \[
          -\rho^{2}\sin\phi\cos\phi(\sin^{2}\theta+\cos^{2}\theta)
            =-\rho^{2}\sin\phi\cos\phi .
        \]
        Multiplied by $\cos\phi$ this is
        $-\rho^{2}\sin\phi\cos^{2}\phi$. Second bracket:
        $\rho\sin^{2}\phi(\cos^{2}\theta+\sin^{2}\theta)=\rho\sin^{2}\phi$.
        Multiplied by $-\rho\sin\phi$: $-\rho^{2}\sin^{3}\phi$. Sum:
        $-\rho^{2}\sin\phi(\cos^{2}\phi+\sin^{2}\phi)=-\rho^{2}\sin\phi$.
        Thus $|\det J|=\rho^{2}\sin\phi$, the spherical volume element.

  \item $U_{x}=a$, $U_{y}=b$. $\dd{y}{x}=-a/b$. \textbf{MRS} $=a/b$
        (constant, perfect substitutes).

  \item Simplify: $f=\dfrac{(x+y)(x^{2}-xy+y^{2})}{x+y}=x^{2}-xy+y^{2}$.
        Then $f_{x}=2x-y$, $f_{y}=-x+2y$.
        \[
          \mathrm{LHS}=x(2x-y)+y(-x+2y)=2x^{2}-2xy+2y^{2} .
        \]
        \[
          \mathrm{RHS}=2f=2(x^{2}-xy+y^{2})=2x^{2}-2xy+2y^{2} .
        \]
        $\mathrm{LHS}=\mathrm{RHS}$. \checkmark

  \item $F=x^{2}+y^{2}+z^{2}-R^{2}$. $F_{x}=2x$, $F_{y}=2y$, $F_{z}=2z$.
        Provided $z\neq 0$:
        \[
          \pd{z}{x}=-\frac{F_{x}}{F_{z}}=-\frac{x}{z},\qquad
          \pd{z}{y}=-\frac{F_{y}}{F_{z}}=-\frac{y}{z} .
        \]
        (Geometrically, the outward normal on the sphere has direction
        $(x,y,z)/R$.)

  \item $f(tx,ty)=(tx)\mathrm{e}^{(ty)/(tx)}=tx\,\mathrm{e}^{y/x}
        =t\,f(x,y)$. \textbf{Degree} $k=1$.

  \item
        \[
          J=\begin{pmatrix} 2u & -2v\\ 2v & 2u\end{pmatrix},\qquad
          \det J=4u^{2}+4v^{2}=4(u^{2}+v^{2}) .
        \]
        Vanishes only at the origin. (This is the complex squaring
        $w=(u+iv)^{2}$.)

  \item $U_{x}=\alpha/x$, $U_{y}=\beta/y$. $U_{x}/U_{y}=(\alpha/\beta)(y/x)$.
        $\dd{y}{x}=-(\alpha/\beta)(y/x)$. \textbf{MRS}
        $=(\alpha/\beta)(y/x)$.

  \item $f=(x^{2}+y^{2})^{1/2}$. $f_{x}=x/\sqrt{x^{2}+y^{2}}$,
        $f_{y}=y/\sqrt{x^{2}+y^{2}}$.
        \[
          \mathrm{LHS}=\frac{x^{2}}{\sqrt{x^{2}+y^{2}}}
            +\frac{y^{2}}{\sqrt{x^{2}+y^{2}}}
            =\frac{x^{2}+y^{2}}{\sqrt{x^{2}+y^{2}}}=\sqrt{x^{2}+y^{2}} .
        \]
        \[
          \mathrm{RHS}=1\cdot f=\sqrt{x^{2}+y^{2}} .
        \]
        $\mathrm{LHS}=\mathrm{RHS}$. \checkmark

  \item $F=\sin(xy)+\cos(x+y)-1/2$. $F_{x}=y\cos(xy)-\sin(x+y)$,
        $F_{y}=x\cos(xy)-\sin(x+y)$.
        \[
          \dd{y}{x} = -\frac{y\cos(xy)-\sin(x+y)}{x\cos(xy)-\sin(x+y)}
            = \frac{\sin(x+y)-y\cos(xy)}{x\cos(xy)-\sin(x+y)} .
        \]

  \item $f(tx,ty)=t^{2}x^{2}+t^{3}y^{3}$. The two terms scale with
        different powers of $t$, so they cannot be factored as
        $t^{k}(x^{2}+y^{3})$. \textbf{Not homogeneous.}

  \item
        \[
          J=\begin{pmatrix} \cos\theta & -r\sin\theta\\
                             \sin\theta & r\cos\theta\end{pmatrix},\qquad
          \det J=r\cos^{2}\theta+r\sin^{2}\theta=r .
        \]
        Hence $|\det J|=r$, and the area element is
        $\mathrm{d}A=r\,\mathrm{d}r\,\mathrm{d}\theta$, the standard
        polar substitution.

  \item Let $S=K^{\rho}+L^{\rho}$. Then $Q=S^{1/\rho}$,
        $Q_{K}=S^{1/\rho-1}K^{\rho-1}$, $Q_{L}=S^{1/\rho-1}L^{\rho-1}$.
        \[
          \frac{Q_{K}}{Q_{L}}=\Bigl(\frac{K}{L}\Bigr)^{\rho-1}
            =\Bigl(\frac{L}{K}\Bigr)^{1-\rho} .
        \]
        $\dd{L}{K}=-(L/K)^{1-\rho}$. \textbf{MRTS} $=(L/K)^{1-\rho}$.

  \item On the smooth piece $\{x<y\}$: $f=x$, $f_{x}=1$, $f_{y}=0$.
        \[
          \mathrm{LHS}=x\cdot 1+y\cdot 0=x=f,\qquad \mathrm{RHS}=1\cdot f=x .
        \]
        $\mathrm{LHS}=\mathrm{RHS}$. \checkmark On $\{x>y\}$ by symmetry
        $\mathrm{LHS}=y=f=\mathrm{RHS}$. \checkmark Euler's identity
        holds wherever the function is differentiable. (At the ridge
        $x=y$ the function is not differentiable.)

  \item By IFT, $\pd{y}{x_{i}}=-F_{x_{i}}/F_{y}$ with
        $F_{y}=x_{2}+3y^{2}$:
        \[
          \pd{y}{x_{1}}=-\frac{2x_{1}}{x_{2}+3y^{2}},\qquad
          \pd{y}{x_{2}}=-\frac{y}{x_{2}+3y^{2}} .
        \]

  \item Numerator has degree $3$, denominator degree $3$. Concretely
        \[
          f(tx,ty,tz)=\frac{t^{3}(x^{2}y+y^{2}z+z^{2}x)}{t^{3}xyz}=f(x,y,z) .
        \]
        \textbf{Degree} $k=0$.

  \item $J=\begin{pmatrix} 1 & \cos y\\ \cos x & 1\end{pmatrix}$,
        $\det J=1-\cos x\cos y$. Vanishes when $\cos x\cos y=1$, i.e.,
        $\cos x=\cos y=1$ or $\cos x=\cos y=-1$, equivalently
        $(x,y)\in 2\pi\mathbb{Z}\times 2\pi\mathbb{Z}$ or
        $(2\mathbb{Z}+1)\pi\times(2\mathbb{Z}+1)\pi$.

  \item Let $S=x^{\rho}+y^{\rho}$, $U=S^{1/\rho}$.
        $U_{x}=S^{1/\rho-1}x^{\rho-1}$, $U_{y}=S^{1/\rho-1}y^{\rho-1}$.
        \[
          \frac{U_{x}}{U_{y}}=\Bigl(\frac{x}{y}\Bigr)^{\rho-1}
            =\Bigl(\frac{y}{x}\Bigr)^{1-\rho} .
        \]
        $\dd{y}{x}=-(y/x)^{1-\rho}$. \textbf{MRS} $=(y/x)^{1-\rho}$.

  \item $f_{x}=y+z$, $f_{y}=x+z$, $f_{z}=x+y$.
        \[
          \mathrm{LHS}=x(y+z)+y(x+z)+z(x+y)=2(xy+yz+zx)=2f .
        \]
        \[
          \mathrm{RHS}=2f .
        \]
        $\mathrm{LHS}=\mathrm{RHS}$. \checkmark

  \item $F=x^{2}+y^{2}-1$. $F_{x}=2x$, $F_{y}=2y$. $\dd{y}{x}=-x/y$.
        At $(\sqrt{3}/2,1/2)$: $\dd{y}{x}=-(\sqrt{3}/2)/(1/2)=-\sqrt{3}$.
        \checkmark

  \item $f(tx,ty)=\bigl((tx)^{1/2}+(ty)^{1/2}\bigr)^{2}
        =(t^{1/2})^{2}(\sqrt{x}+\sqrt{y})^{2}=t f(x,y)$. \textbf{Degree}
        $k=1$. For Euler's theorem, write $A=\sqrt{x}+\sqrt{y}$,
        $f=A^{2}$. Then $f_{x}=2A\cdot\dfrac{1}{2\sqrt{x}}=A/\sqrt{x}$,
        $f_{y}=A/\sqrt{y}$.
        \[
          \mathrm{LHS}=x\cdot\frac{A}{\sqrt{x}}+y\cdot\frac{A}{\sqrt{y}}
            =A(\sqrt{x}+\sqrt{y})=A^{2}=f .
        \]
        \[
          \mathrm{RHS}=1\cdot f=A^{2} .
        \]
        $\mathrm{LHS}=\mathrm{RHS}$. \checkmark

  \item
        \[
          J=\begin{pmatrix} 1/y & -x/y^{2}\\ -y/x^{2} & 1/x\end{pmatrix},
          \qquad \det J=\frac{1}{xy}-\frac{xy}{x^{2}y^{2}}
            =\frac{1}{xy}-\frac{1}{xy}=0 .
        \]
        The Jacobian is identically singular: because
        $(x/y)\cdot(y/x)=1$, the image lies on the curve $f_{1}f_{2}=1$
        (one-dimensional), so the map is not locally invertible
        anywhere.

  \item From $\ln Q=\alpha_{0}+\alpha_{K}\ln K+\alpha_{L}\ln L$, we get
        $Q=AK^{\alpha_{K}}L^{\alpha_{L}}$ with $A=\mathrm{e}^{\alpha_{0}}$
        (Cobb--Douglas). Then $Q_{K}=A\alpha_{K}K^{\alpha_{K}-1}L^{\alpha_{L}}$,
        $Q_{L}=A\alpha_{L}K^{\alpha_{K}}L^{\alpha_{L}-1}$.
        \[
          \frac{Q_{K}}{Q_{L}}=\frac{\alpha_{K}}{\alpha_{L}}\cdot\frac{L}{K} .
        \]
        $\dd{L}{K}=-(\alpha_{K}/\alpha_{L})(L/K)$. \textbf{MRTS}
        $=(\alpha_{K}/\alpha_{L})(L/K)$.

  \item $Q=K^{1/3}L^{1/3}M^{1/3}$ is homogeneous of degree
        $1/3+1/3+1/3=1$. $Q_{K}=\tfrac13K^{-2/3}L^{1/3}M^{1/3}=Q/(3K)$,
        and similarly $Q_{L}=Q/(3L)$, $Q_{M}=Q/(3M)$.
        \[
          \mathrm{LHS}=K\cdot\frac{Q}{3K}+L\cdot\frac{Q}{3L}
            +M\cdot\frac{Q}{3M}=\frac{Q}{3}+\frac{Q}{3}+\frac{Q}{3}=Q .
        \]
        \[
          \mathrm{RHS}=1\cdot Q=Q .
        \]
        $\mathrm{LHS}=\mathrm{RHS}$. \checkmark

  \item Write $f_{1}=x+y+z$, $f_{2}=xy+yz+zx$, $f_{3}=xyz$ (the
        elementary symmetric polynomials $e_{1},e_{2},e_{3}$). Then
        \[
          J=\begin{pmatrix}
            1 & 1 & 1\\
            y+z & x+z & x+y\\
            yz & xz & xy
          \end{pmatrix} .
        \]
        For the determinant, subtract column $1$ from columns $2$ and
        $3$ to get
        \[
          \det J = \det\begin{pmatrix}
            1 & 0 & 0\\
            y+z & x-y & x-z\\
            yz & z(x-y) & y(x-z)
          \end{pmatrix}
          = 1\cdot\det\begin{pmatrix}
            x-y & x-z\\ z(x-y) & y(x-z)
          \end{pmatrix}
        \]
        \[
          = (x-y)(x-z)\det\begin{pmatrix} 1 & 1\\ z & y\end{pmatrix}
          = (x-y)(x-z)(y-z) .
        \]
        Thus $\det J=-(x-y)(y-z)(z-x)$ (the Vandermonde discriminant, up
        to sign), which vanishes exactly when two of $x,y,z$ coincide:
        the symmetric-polynomial map fails to be locally invertible
        at points with repeated coordinates, consistent with the fact
        that an unordered triple $\{x,y,z\}$ has a unique sorted
        representative.
  \item $f'(x)=5x^{4}+2\ge2>0$ for all $x$, so $f$ is strictly increasing,
        hence injective, and the theorem applies everywhere. Solve
        $x^{5}+2x+1=4$: $x=1$ works ($1+2+1=4$), and monotonicity makes it
        the only root. So
        $\bigl(f^{-1}\bigr)'(4)=1/f'(1)=1/7$.
  \item $f(1)=1+3-2=2$, and $f'(x)=3x^{2}+3$, so $f'(1)=6$ and
        $\bigl(f^{-1}\bigr)'(2)=\tfrac16$. ($f'>0$ everywhere, so the
        inverse is global.)
  \item $f(0)=1+0=1$, and $f'(x)=\mathrm{e}^{x}+1$, so $f'(0)=2$ and
        $\bigl(f^{-1}\bigr)'(1)=\tfrac12$.
  \item $f(1)=1+\ln1=1$, and $f'(x)=1+1/x$, so $f'(1)=2$ and
        $\bigl(f^{-1}\bigr)'(1)=\tfrac12$.
  \item Differentiate both sides with respect to $y$. The left side is a
        composite, so the chain rule gives
        $f'\bigl(f^{-1}(y)\bigr)\cdot\bigl(f^{-1}\bigr)'(y)$; the right side
        gives $1$. Hence
        $\bigl(f^{-1}\bigr)'(y)=1/f'\bigl(f^{-1}(y)\bigr)$, the division
        being legitimate exactly when $f'\neq0$ there.
  \item Invertibility survives: $x^{3}$ is strictly increasing and onto, so
        $f^{-1}(y)=y^{1/3}$ exists globally. What fails is
        \emph{differentiability of the inverse at $0$}:
        $\bigl(y^{1/3}\bigr)'=\tfrac13y^{-2/3}\to\infty$ as $y\to0$. The
        hypothesis $f'\neq0$ is sufficient for a differentiable local
        inverse, not necessary for an inverse.
  \item $f'(x)=2x\neq0$ for $x\neq0$, so the theorem applies at every point
        except the origin. Around $x=2$ the largest interval of
        injectivity is $(0,\infty)$: on it $f$ is strictly increasing, while
        any interval containing $0$ in its interior contains a pair
        $\pm a$ with equal images.
  \item On $(-\pi/2,\pi/2)$ the sine is strictly increasing, and
        $\pi/3$ lies inside. By the theorem, with
        $\sin(\pi/3)=\sqrt3/2$ and $\cos(\pi/3)=\tfrac12$,
        \[
          \bigl(\sin^{-1}\bigr)'\Bigl(\tfrac{\sqrt3}{2}\Bigr)
          =\frac{1}{\cos(\pi/3)}=2 .
        \]
        Directly: $\arcsin'(y)=1/\sqrt{1-y^{2}}
        =1/\sqrt{1-\tfrac34}=1/\tfrac12=2$. \checkmark
  \item $J=\begin{pmatrix}1&1\\1&-1\end{pmatrix}$, $\det J=-2\neq0$
        everywhere. Solving $u=x+y$, $v=x-y$ gives
        $x=\tfrac{u+v}{2}$, $y=\tfrac{u-v}{2}$, so
        $J_{\mathbf{f}^{-1}}
        =\begin{pmatrix}1/2&1/2\\1/2&-1/2\end{pmatrix}$, which is indeed
        $J^{-1}=-\tfrac12\begin{pmatrix}-1&-1\\-1&1\end{pmatrix}$.
        \checkmark\ The map is linear with $\det\neq0$, so here local and
        global invertibility coincide.
  \item $J=\begin{pmatrix}2x&-2y\\2y&2x\end{pmatrix}$ and
        $\det J=4(x^{2}+y^{2})$, so at $(1,2)$ it is $4\cdot5=20\neq0$: a
        differentiable local inverse exists near
        $\mathbf{f}(1,2)=(-3,4)$.
  \item $\mathbf{f}(-x,-y)=\bigl((-x)^{2}-(-y)^{2},\,2(-x)(-y)\bigr)
        =(x^{2}-y^{2},\,2xy)=\mathbf{f}(x,y)$, so every point and its
        negative share an image --- for instance
        $\mathbf{f}(1,1)=\mathbf{f}(-1,-1)=(0,2)$. The map is two-to-one,
        and no global inverse exists. Local invertibility at every point
        does not assemble into global invertibility.
  \item $J=\mathrm{e}^{x}\begin{pmatrix}\cos y&-\sin y\\
        \sin y&\cos y\end{pmatrix}$, so
        $\det J=\mathrm{e}^{2x}(\cos^{2}y+\sin^{2}y)=\mathrm{e}^{2x}>0$
        everywhere. Yet $\mathbf{f}(x,y+2\pi)=\mathbf{f}(x,y)$. There is no
        contradiction: the theorem promises an inverse on \emph{some
        neighbourhood} of each point, and those neighbourhoods (width less
        than $2\pi$ in $y$) never join up into one.
  \item $J=\begin{pmatrix}\cos\theta&-r\sin\theta\\
        \sin\theta&r\cos\theta\end{pmatrix}$, so $\det J=r$. The theorem
        says nothing on the line $r=0$ --- the origin, where every $\theta$
        gives the same point and the map is infinitely many-to-one.
  \item $J=\begin{pmatrix}1&3y^{2}\\3x^{2}&1\end{pmatrix}$, which at
        $(0,0)$ is the identity, with $\det=1\neq0$. Hence a local inverse
        exists near $\mathbf{f}(0,0)=(0,0)$ and
        $J_{\mathbf{f}^{-1}}(0,0)=I^{-1}=I$.
  \item $J=\begin{pmatrix}1&0&0\\1&1&0\\1&1&1\end{pmatrix}$, lower
        triangular, so $\det J=1\neq0$. From $u=x$, $v=x+y$, $w=x+y+z$:
        $x=u$, $y=v-u$, $z=w-v$, giving
        $J_{\mathbf{f}^{-1}}=\begin{pmatrix}1&0&0\\-1&1&0\\
        0&-1&1\end{pmatrix}$. Multiplying the two matrices returns $I$.
        \checkmark
  \item It permits: near $\mathbf{p}_{0}$ the system can in principle be
        solved for prices in terms of quantities, the inverse demand
        $\mathbf{p}(\mathbf{q})$ is differentiable, and its derivative
        matrix is $\bigl[J_{\mathbf{f}}\bigr]^{-1}$ --- so the price
        responses can be computed from the quantity responses without ever
        solving the system. It does not permit: any claim that the same
        quantities cannot arise from a different price vector elsewhere.
        Global identification is a separate question, and $\det J\neq0$ is
        no evidence for it.
\end{solutions}

\section{Additional Problems}

\begin{problems}
  \item Use IFT to find $\dd{y}{x}$ for $x^{2}+2xy+y^{2}=9$.
  \item Verify homogeneity and find the degree of $f(x,y)=x^{3}y+xy^{3}$.
  \item Compute the Jacobian and its determinant for
        $\mathbf{f}(x,y)=(2x+y,\;x-y)$.
  \item For $U(x,y)=x^{2}+y$, find the signed slope and unsigned MRS
        using IFT.
  \item Verify Euler's theorem for $f(x,y)=x^{3}+2x^{2}y-y^{3}$
        (degree $3$).
  \item Use IFT to find $\dd{y}{x}$ for $x^{3}+y^{3}+xy=5$.
  \item Verify homogeneity and find the degree of
        $f(x,y)=\dfrac{x^{4}+y^{4}}{x^{2}+y^{2}}$.
  \item Compute the Jacobian and its determinant for
        $\mathbf{f}(x,y)=(x^{2}y,\;xy^{2})$.
  \item For $U(x,y)=x^{1/3}y^{2/3}$, find the signed slope and MRS via
        IFT.
  \item Verify Euler's theorem for $Q(K,L)=K^{1/2}L^{1/2}$ (degree $1$).
  \item Use IFT to find $\dd{y}{x}$ for $\ln(xy)+x+y=4$.
  \item Decide whether $f(x,y)=x+y^{2}$ is homogeneous, with
        justification.
  \item Compute the Jacobian determinant of
        $\mathbf{f}(x,y,z)=(x+y,\;y+z,\;x+z)$.
  \item For $Q(K,L)=aK+bL$ ($a,b>0$), find the MRTS.
  \item Verify Euler's theorem for $f(x,y,z)=xy+yz+xz$ (degree $2$,
        three-variable).
  \item For $f(x)=x^{7}+x^{3}+x$, compute $\bigl(f^{-1}\bigr)'(3)$.
  \item For $f(x)=2x+\sin x$, compute $\bigl(f^{-1}\bigr)'(0)$.
  \item For $f(x)=\ln x+x^{2}$ on $x>0$, compute $\bigl(f^{-1}\bigr)'(1)$.
  \item For $f(x)=\tan x$ on $(-\pi/2,\pi/2)$, compute
        $\bigl(f^{-1}\bigr)'(1)$ and check it against $\arctan'$.
  \item For $\mathbf{f}(x,y)=(2x+3y,\;x+2y)$, give $\det J$ and
        $J_{\mathbf{f}^{-1}}$, and say why the inverse is global here.
  \item For $\mathbf{f}(x,y)=(x^{3},\;y^{3})$: where does $\det J$ vanish,
        and is the map invertible there anyway?
  \item For $\mathbf{f}(x,y)=(x+\mathrm{e}^{y},\;y+\mathrm{e}^{x})$,
        evaluate $\det J$ at the origin and say what the theorem concludes.
  \item Let $\mathbf{f}(\mathbf{x})=A\mathbf{x}$ with $A$ an $n\times n$
        matrix. Show that $\det J$ is constant, and that local and global
        invertibility coincide.
  \item For $\mathbf{f}(x,y)=(x^{2}+y^{2},\;x^{2}-y^{2})$ on $x,y>0$,
        compute $\det J$ and decide local invertibility.
  \item State in one sentence each what $f'(x_{0})\neq0$ and
        $\det J(\mathbf{x}_{0})\neq0$ do and do not guarantee.
\end{problems}

\section{Additional Solutions}

\begin{solutions}
  \item $F=x^{2}+2xy+y^{2}-9=(x+y)^{2}-9$. $F_{x}=2(x+y)$,
        $F_{y}=2(x+y)$.
        \[
          \dd{y}{x}=-\frac{F_{x}}{F_{y}}=-1 .
        \]
        (The curve is two lines $x+y=\pm 3$, each with slope $-1$.)

  \item $f(tx,ty)=(tx)^{3}(ty)+(tx)(ty)^{3}=t^{4}(x^{3}y+xy^{3})
        =t^{4}f(x,y)$. \textbf{Degree $k=4$.}

  \item
        \[
          J=\begin{pmatrix} 2 & 1\\ 1 & -1\end{pmatrix},\qquad
          \det J=-2-1=-3 .
        \]
        Globally invertible; areas scale by $|{-3}|=3$ (with orientation
        reversal).

  \item $U_{x}=2x$, $U_{y}=1$. By IFT $\dd{y}{x}=-2x$. \textbf{MRS}
        $=2x$ (unsigned).

  \item $f_{x}=3x^{2}+4xy$, $f_{y}=2x^{2}-3y^{2}$.
        \[
          \mathrm{LHS}=x(3x^{2}+4xy)+y(2x^{2}-3y^{2})
          =3x^{3}+4x^{2}y+2x^{2}y-3y^{3}=3x^{3}+6x^{2}y-3y^{3} .
        \]
        \[
          \mathrm{RHS}=3f=3(x^{3}+2x^{2}y-y^{3})
          =3x^{3}+6x^{2}y-3y^{3} .
        \]
        $\mathrm{LHS}=\mathrm{RHS}$. \checkmark

  \item $F=x^{3}+y^{3}+xy-5$. $F_{x}=3x^{2}+y$, $F_{y}=3y^{2}+x$.
        \[
          \dd{y}{x}=-\frac{3x^{2}+y}{3y^{2}+x} .
        \]

  \item $f(tx,ty)=\dfrac{t^{4}(x^{4}+y^{4})}{t^{2}(x^{2}+y^{2})}
        =t^{2}f(x,y)$. \textbf{Degree $k=2$.}

  \item
        \[
          J=\begin{pmatrix} 2xy & x^{2}\\ y^{2} & 2xy\end{pmatrix},\qquad
          \det J=4x^{2}y^{2}-x^{2}y^{2}=3x^{2}y^{2} .
        \]
        Singular when $x=0$ or $y=0$.

  \item $U_{x}=\tfrac13 x^{-2/3}y^{2/3}$,
        $U_{y}=\tfrac23 x^{1/3}y^{-1/3}$.
        $U_{x}/U_{y}=\tfrac12\cdot y/x$. By IFT $\dd{y}{x}=-y/(2x)$.
        \textbf{MRS} $=y/(2x)$.

  \item $Q_{K}=\tfrac12 K^{-1/2}L^{1/2}$,
        $Q_{L}=\tfrac12 K^{1/2}L^{-1/2}$.
        \[
          \mathrm{LHS}=K\cdot\tfrac12 K^{-1/2}L^{1/2}
          +L\cdot\tfrac12 K^{1/2}L^{-1/2}
          =\tfrac12 K^{1/2}L^{1/2}+\tfrac12 K^{1/2}L^{1/2}
          =K^{1/2}L^{1/2} .
        \]
        \[
          \mathrm{RHS}=1\cdot Q=K^{1/2}L^{1/2} .
        \]
        $\mathrm{LHS}=\mathrm{RHS}$. \checkmark

  \item $F=\ln(xy)+x+y-4=\ln x+\ln y+x+y-4$. $F_{x}=1/x+1$,
        $F_{y}=1/y+1$.
        \[
          \dd{y}{x}=-\frac{1/x+1}{1/y+1}=-\frac{y(1+x)}{x(1+y)} .
        \]

  \item $f(tx,ty)=tx+t^{2}y^{2}$. The two terms scale with different
        powers of $t$, so they cannot share a common $t^{k}$ factor.
        \textbf{Not homogeneous.}

  \item
        \[
          J=\begin{pmatrix}
            1 & 1 & 0\\
            0 & 1 & 1\\
            1 & 0 & 1
          \end{pmatrix} .
        \]
        Expanding along the first row:
        $\det J=1\cdot(1\cdot 1-1\cdot 0)-1\cdot(0\cdot 1-1\cdot 1)+0
        =1-(-1)+0=2$.

  \item Linear: $Q_{K}=a$, $Q_{L}=b$. $\dd{L}{K}=-a/b$. \textbf{MRTS}
        $=a/b$ (constant, perfect substitutes).

  \item $f_{x}=y+z$, $f_{y}=x+z$, $f_{z}=x+y$.
        \[
          \mathrm{LHS}=x(y+z)+y(x+z)+z(x+y)=2(xy+yz+xz)=2f .
        \]
        \[
          \mathrm{RHS}=2f .
        \]
        $\mathrm{LHS}=\mathrm{RHS}$. \checkmark
  \item $f(1)=1+1+1=3$ and $f'(x)=7x^{6}+3x^{2}+1$, so $f'(1)=11$ and
        $\bigl(f^{-1}\bigr)'(3)=\tfrac{1}{11}$. ($f'>0$ everywhere, so the
        pre-image is unique.)
  \item $f(0)=0$ and $f'(x)=2+\cos x\ge1>0$, so $f'(0)=3$ and
        $\bigl(f^{-1}\bigr)'(0)=\tfrac13$.
  \item $f(1)=\ln1+1=1$ and $f'(x)=1/x+2x$, so $f'(1)=3$ and
        $\bigl(f^{-1}\bigr)'(1)=\tfrac13$.
  \item $\tan(\pi/4)=1$ and $f'(x)=1/\cos^{2}x$, so
        $f'(\pi/4)=1/(\sqrt2/2)^{2}=2$ and
        $\bigl(f^{-1}\bigr)'(1)=\tfrac12$. Directly:
        $\arctan'(1)=1/(1+1^{2})=\tfrac12$. \checkmark
  \item $J=\begin{pmatrix}2&3\\1&2\end{pmatrix}$ is constant with
        $\det J=4-3=1\neq0$, and
        $J_{\mathbf{f}^{-1}}=J^{-1}
        =\begin{pmatrix}2&-3\\-1&2\end{pmatrix}$. The map is linear, so the
        local inverse supplied by the theorem is the global one:
        $\mathbf{x}=A^{-1}\mathbf{y}$ is defined on all of $\R^{2}$.
  \item $J=\begin{pmatrix}3x^{2}&0\\0&3y^{2}\end{pmatrix}$ and
        $\det J=9x^{2}y^{2}$, which vanishes on both axes. The map is
        nevertheless a bijection of $\R^{2}$, with inverse
        $(u,v)\mapsto(u^{1/3},v^{1/3})$ --- but that inverse is not
        differentiable where $\det J=0$. Same lesson as $x^{3}$ in one
        variable.
  \item $J=\begin{pmatrix}1&\mathrm{e}^{y}\\ \mathrm{e}^{x}&1\end{pmatrix}$,
        so at $(0,0)$ it is
        $\begin{pmatrix}1&1\\1&1\end{pmatrix}$ with $\det J=0$. The theorem
        concludes \textbf{nothing}: its hypothesis fails, and that is not
        the same as showing no local inverse exists. A separate argument
        would be needed either way.
  \item Each component of $A\mathbf{x}$ is linear, so
        $\partial(A\mathbf{x})_{i}/\partial x_{j}=a_{ij}$ and $J=A$ at every
        point --- constant, with $\det J=\det A$. If $\det A\neq0$ then
        $\mathbf{x}=A^{-1}\mathbf{y}$ inverts the map on all of $\R^{n}$; if
        $\det A=0$ then $A$ has a non-trivial null space and the map is not
        injective anywhere. So for linear maps the local condition is
        already global.
  \item $J=\begin{pmatrix}2x&2y\\2x&-2y\end{pmatrix}$ and
        $\det J=-4xy-4xy=-8xy$, which is non-zero throughout $x,y>0$. So a
        differentiable local inverse exists at every point of the open
        quadrant.
  \item $f'(x_{0})\neq0$ guarantees an interval around $x_{0}$ on which $f$
        is invertible with a differentiable inverse of derivative
        $1/f'(x_{0})$; it does not guarantee that $f$ is injective anywhere
        beyond that interval. $\det J(\mathbf{x}_{0})\neq0$ guarantees the
        same locally in $n$ dimensions, with the inverse's Jacobian equal to
        $\bigl[J(\mathbf{x}_{0})\bigr]^{-1}$; it likewise says nothing about
        global injectivity.
\end{solutions}
